\chapter{Literature Study}
\label{chap:literature}

This chapter presents the literature study conducted prior to the research described in this thesis. The published literature on pipeline leak detection is extensive and covers a wide range of approaches and problem formulations. Much of the existing research focuses on leak detection in pipe networks consisting of multiple interconnected pipes, whereas a smaller body of research addresses the detection of multiple leaks within a single pipeline. For completeness, \autoref{sec:lit:pipe_network} provides a brief introduction to research on leak detection in pipe networks.\\

Some methods only support sequential leak detection. In these approaches, monitoring begins under the assumption that the pipeline is healthy and free of defects. It is further assumed that leaks do not occur simultaneously. Multiple leaks can then be treated as a sequence of single leak detection problems, where the effect of each detected leak is accounted for before the next leak is identified. An example of this approach is presented by \cite{verde2005}. Such methods fall outside the scope of this research, and the literature review will be limited to methods that address simultaneous leaks.\\

The detection of one or two leaks in a single pipeline using only endpoint sensors is well documented in the literature, whereas methods that address more than two concurrent leaks are rare. The reviewed research can be broadly classified according to whether the available measurements are time-varying transient or steady-state. \autoref{sec:lit:transients} presents an overview of relevant literature based on transient measurements, and \autoref{sec:lit:steadystate} reviews research based on steady-state measurements, which relate closest to the work in this thesis. This section further discusses whether existing methods for two leak detection are feasible and practical to extend to an arbitrary number of leaks.

\newpage
\section{Leak Identification in Pipe Networks}
\label{sec:lit:pipe_network}

Leak identification in pipe networks is commonly formulated as an inverse problem. The early formulation in \cite{pudar1992} estimates equivalent leak-orifice areas at candidate network nodes using pressure and flow measurements. The method shows that the result depends strongly on the number and placement of sensors, as well as on uncertainties in pipe-friction parameters. Inverse transient analysis was extended to pipe networks in \cite{liggett1994}. The publication argues that transient pressure measurements provide more information than steady-state measurements, enabling the simultaneous estimation of leak and pipe-friction parameters. Both methods represent the water distribution network as a graph, where potential leak locations are modeled as nodes, rather than as continuous positions along a pipe.\\

More recent work in \cite{molno2024conditions} studies leak localization in potential-flow networks, where leaks can be located anywhere along edges, rather than only at predefined nodes. The work introduces “one-way” edges between sensor nodes where the potential difference varies monotonically with an assumed leak position. Based on these edges, conditions on the network's graph structure and potential sensor placement are derived, under which it is guaranteed that the leak can be limited to a small set of locations.\\

A specific case of a network of parallel pipes was studied in \cite{molno2024parallel}. One steady-state measurement can yield a single leak equivalent for each pipe. Provided the pipes have different head-loss characteristics, the paper shows that only the true leak will have a constant single leak equivalent across multiple steady-states. Hence, measurements from different operating states can be combined to identify the leaking pipe and the leak's location along that pipe.\\

Modern data-driven methods approach network-level leak localization by learning how leaks change pressure and flow patterns across the network. A representative early deep-learning contribution is presented in \cite{hu2021multiscale}. Instead of using each pipe as a separate classification label, the network is first divided into leak zones using DBSCAN clustering, which groups pipes with similar hydraulic behavior. A multiscale fully convolutional network is then trained on simulated pressure and flow residuals to identify the most likely zone. This makes leak localization more practical for large networks, but limits the result to a leak area rather than an actual position along a pipe.\\

A generative approach that reduces the need for labeled leak data is presented in \cite{rajabi2023}. The method uses a conditional generative adversarial network to learn the normal relation between demand and pressure distributions. The model is trained to translate a demand pattern into the pressure pattern expected in a leak-free network. A leak is detected when the measured pressure deviates from the expected pattern, and localization is based on the location where local similarity is lowest. The method, therefore, uses the spatial pressure pattern of the entire network rather than manually selected residual features. However, it still depends on a hydraulic model to generate representative leak-free training data, so its performance remains tied to the quality of that model.\\

A more recent graph-based method is presented in \cite{zhang2026aignn}. Here, the water distribution system is treated directly as a graph, and the neural network is first trained to imitate a classical max-flow algorithm. This is intended to enable the model to learn how information flows through connected pipes and nodes before it is applied to leak localization. One network reconstructs the current pressure field from available measurements, while another predicts the expected pressure field from earlier measurements. A leak is indicated by discrepancies between these two fields. The main contribution is improved generalization, in which the model is encouraged to follow the network structure rather than "memorize" simulated leak examples. Like the other data-driven methods, it remains a network-level localization method, not a method for estimating an exact leak location along a pipe.

\section{Leak Identification in Pipeline with Transient Measurements}
\label{sec:lit:transients}

Transient methods use pressure and flow measurements collected while the flow is changing. This additional information is useful for concurrent leaks, since a single steady-state may not uniquely identify the leak configuration \cite{verde2004}.

\subsection*{Observer and Kalman Filter Methods}

Observer-based methods estimate leak locations and sizes, as well as the pipe's hydraulic states. In \cite{besancon2007}, a specified number of leaks is estimated from endpoint measurements and transient excitation. The method is demonstrated for two simultaneous leaks. Discretized pipe models are used in \cite{torres2009,verde2001}. Instead of allowing leaks at any position, the pipe is divided into candidate locations or sections. Multiple leaks are identified by identifying which candidates act as leak points. This reduces the location problem to a finite set of choices, but the accuracy of the location depends on how finely the pipe is divided.

\subsection*{Inverse Transient Analysis}

Inverse transient analysis identifies leaks by comparing measured transient signals with simulated signals from a pipe model. A two-leak formulation in \cite{diaz2021} indentifies leaks from inlet and outlet measurements. Other approaches first reduce the search area before applying inverse transient analysis. In \cite{nadian2025}, wavelet analysis is used to find approximate leak positions from reflected pressure waves. In \cite{cherif2025}, matched field processing defines a smaller search interval before particle swarm optimization is applied. Both approaches focus on two simultaneous leaks.\\

Some inverse transient methods compare pressure signals in the frequency domain. In \cite{hyun2018}, the leak locations and sizes are estimated by matching the measured frequency response. The formulation in \cite{wang2018} represents the measured response as a sum of contributions from several leaks. This makes it possible to estimate multiple leak locations simultaneously. The extension in \cite{wang2019} estimates one leak contribution at a time and uses model selection criteria to estimate the number of leaks.

\subsection*{Equivalent-Leak Based Transient Methods}

Equivalent leak-based transient methods form a small line of work built around the same observation. The analysis in \cite{verde2004} shows that two leaks can exhibit the same steady-state boundary behavior as a single virtual leak, so steady-state data alone do not isolate the real leak pair. This virtual leak is then used to define a smaller family of possible two-leak models. The method in \cite{verde2007} combines this family with a spatially discretized transient model. Steady-state data select the relevant family, and transient endpoint data are used to choose the pair of grid locations that yields the smallest output error. The method is extended in \cite{verde2014}, which introduces a parameterized transient model for the same two-leak problem. This reduces the search from four unknown leak parameters to two before matching measured endpoint pressure and flow signals.

\section{Leak Identification in Pipeline with steady-state Measurements}
\label{sec:lit:steadystate}

\subsection{Utilizing Multiple Pressure Sensors}

The study in \cite{ostapkowicz2023} proposes a static model-based approach for identifying two leaks in a pipeline using steady-state flow rate and pressure measurements at the inlet and outlet, complemented by pressure measurements at several intermediate locations along the pipeline. The framework introduces two identification methods.\\

The first method is based on a steady-state hydraulic model of a pipeline with two leaks. The leak locations are treated as decision variables and are identified by minimizing an objective function defined as the squared deviation between measured pressures and model-predicted pressures at the intermediate sensor locations. The objective function can be written as
\begin{equation*}
J(z_{u1}, z_{u2})
= \sum_{i=2}^{N-1} \left( p_i^{\text{meas}} - p_i^{\text{mod}}(z_{u1}, z_{u2}) \right)^2,
\end{equation*}

where $p_i^{\text{meas}}$ denotes the measured pressure at sensor $i$, $p_i^{\text{mod}}$ is the pressure computed from the static model, and $z_{u1}$, $z_{u2}$ are the leak locations. Once the optimal leak positions are obtained, the corresponding leak magnitudes are computed from steady-state flow-balance relations among the inlet, outlet, and intermediate pipeline sections.\\

The second method also employs a static pipeline model but replaces global optimization with a pressure gradient-based procedure. Local pressure gradients between adjacent measurement points are first computed as
\begin{equation*}
g_k = z_{k+1} - z_k \frac{p_k}{p_{k+1}}
\end{equation*}
and gradient indicator functions are used to identify pipeline segments that contain leaks. After the leak-containing segments are isolated, analytical expressions derived from the static pressure-drop relations are applied to estimate both the leak locations and magnitudes within each segment.\\

Both methods were experimentally validated for simultaneous double-leak scenarios. A limitation of the proposed approach is its strong dependence on pressure measurements at multiple intermediate locations along the pipeline. The availability of several internal pressure sensors is essential for both the optimization-based and gradient-based procedures. As a result, the proposed methods are not feasible for pipeline configurations where only inlet and outlet measurements are available.

\subsection{Exploiting Single-Leak Equivalent}
\label{sec:peralta_study}

It is shown in \cite{verde2003} that input-output measurements at a single steady-state operating point are insufficient to detect two or more leaks in a pipeline. The method proposed in \cite{peralta2023} identifies two concurrent leaks in a single pipeline using measurements of flow rate and pressure at the pipe's inlet and outlet, collected at two different steady-state operating points.\\

From a single operating point, a single-leak equivalent position $L_{eq}$ is computed as
\begin{equation*}
L_{eq} =
\frac{h_0 - h_\ell - \ell \, \theta \, q_\ell^2}
{\theta \left( q_0^2 - q_\ell^2 \right)},
\end{equation*}
where $h_0$ and $h_\ell$ denote the inlet and outlet pressure, $q_0$ and $q_\ell$ denote the inlet and outlet flow rates, $\ell$ is the known pipe lenght and $\theta$ is a known pipe parameter.

For a pipe with two leaks located at positions $L_1$ and $L_1 + L_2$, the equivalent single-leak position satisfies the inequality
\begin{equation}
\label{eq:lit_ineq}
0 < L_1 < L_{eq} < L_1 + L_2 < \ell.
\end{equation}
The key idea is that the single-leak equivalent position can be estimated from either the upstream or the downstream leak. For each operating point $i$, this redundant description is given by
\begin{equation}
\label{eq:lit_lf_estimate}
\left.\hat{L}_{eq}\right\vert_i =
\begin{cases}
\left.\hat{L}^u_{eq}\right\vert_i
:=
L_1 + L_2
\left(
\dfrac{q_m^2(L_1,\ c_1) - q_{\ell}^2}
{q_0^2 -  q_{\ell}^2}
\right), \\[6pt]
\left.\hat{L}^{d}_{eq}\right\vert_i
:=
L_1 + L_2
\left(
\dfrac{q_m^2(L_1,L_2,c_2) -  q_{\ell}^2}
{q_0^2 -  q_{\ell}^2}
\right),
\end{cases}
\end{equation}
where $q_m$ is the unknown flow between the two leaks, and $c_1$ and $c_2$ are the leak coefficients.\\

This redundant formulation is used to define a cost function based on the difference between the estimated leak equivalents and the actual equivalent position. The cost function is defined as
\begin{equation*}
J(e^u,e^d) =
\frac{1}{m}
\sum_{i=1}^{m}
\left(
\left.e^u\right\vert_i^2 + \left.e^d\right\vert_i^2
\right),
\end{equation*}
where $m$ is the number of measured operating points. The upstream and downstream estimation errors are defined as
\begin{equation*}
\left.e^u\right\vert_i = \left.L_{eq}\right\vert_i - \left.\hat{L}^{u}_{eq}\right\vert_i,
%\quad \text{and}
\quad \left.e^d\right\vert_i = \left.L_{eq}\right\vert_i - \left.\hat{L}^{d}_{eq}\right\vert_i
\quad \forall i.
\end{equation*}

The cost function is minimized subject to the invariance of the parameters $c_1$ and $c_2$ and the constraint
\begin{equation*}
0 < L_1 < \left.L_{eq}\right\vert_i < L_1 + L_2 < \ell.
\end{equation*}

In this publication, the minimization is performed using a gradient-based update law that adjusts only the upstream leak position,
\begin{equation*}
L_1(k+1) =
L_1(k) -
\alpha
\frac{\partial J(k)}{\partial L_1(k)},
\end{equation*}
with the constraint $0 < L_1 < L_{eq}$. \\

The approach is extended in \cite{peralta2024} by simultaneously optimizing leak locations and reducing estimation errors on the same dataset. The new update law is given by
\begin{equation*}
\begin{bmatrix}
L_{1}(k+1) \\
L_{3}(k+1)
\end{bmatrix}
=
\begin{bmatrix}
L_{1}(k) \\
L_{3}(k)
\end{bmatrix}
+
\begin{bmatrix}
\alpha \, \dfrac{\partial J(k)}{\partial L_1(k)} \\
\beta \, \dfrac{\partial J(k)}{\partial L_3(k)}
\end{bmatrix},
\end{equation*}
with the constraints $0 < L_{1} < {L}_{eq}$ and $0 < L_{3} < \ell - {L}_{eq}$ where $L_3 = \ell- (L_1 + L_2)$. A case study of this method is presented in \cite{peralta2025} where both simulation and experimental results are shown with a relative leak location error of less than 7\%.\\

This method is based on the inequality in \cref{eq:lit_ineq}, where the single-leak equivalent will always lie between the two actual leaks. This allows for the consistent upstream and downstream estimates shown in \cref{eq:lit_lf_estimate}. To extend this method beyond two leaks, it seems necessary that a similar inequality hold, such that the single-leak equivalent has the same number of real leaks upstream and downstream for all steady-state operating points. Otherwise, there will be no consistent algebraic upstream and downstream estimate for the single-leak equivalent across measurements. It can be shown (see \autoref{app:3_leak_eq_proof}) that the single-leak equivalent in a pipe with three leaks can be located both upstream and downstream of the second leak by adjusting the steady-state operating point. Based on this result, it seems difficult to extend the method proposed in \cite{peralta2023} to handle more than two leaks.  

\subsection{Monte Carlo Simulation}

The method proposed in \cite{torres2024} addresses the localization of two simultaneous leaks under steady-state conditions using only inlet and outlet measurements. The approach formulates the inverse problem as a constrained Monte Carlo identification procedure based on algebraic hydraulic models. Instead of enforcing a unique solution, the method characterizes the feasible parameter space through probability distributions.\\

The algorithm consists of four steps:
\begin{enumerate}
    \item \textbf{Equivalent leak estimation:}  
    An equivalent single-leak model is computed from steady-state inlet and outlet pressure and flow rate measurements using a method similar to that presented in \cite{peralta2023}. Similarly, the single-leak equivalent is used as spatial bounds for the real upstream and downstream leaks.

    \item \textbf{Forward Monte Carlo sampling:}  
    Candidate values for the downstream leak location and emitter coefficient are randomly generated from uniform distributions bounded by the pipe geometry and the equivalent leak parameters. For each realization, the corresponding upstream leak location and coefficient are computed analytically using steady-state flow and pressure relations. Realizations that violate spatial ordering or hydraulic feasibility constraints are discarded.

    \item \textbf{Probability density estimation:}  
    The retained samples of the upstream leak parameters are used to fit probability density functions via maximum likelihood estimation. These distributions provide a statistical description of the upstream leak location and magnitude consistent with the measurements and imposed constraints.

    \item \textbf{Backward Monte Carlo refinement:}  
    New samples of the upstream leak parameters are drawn from the estimated probability density functions, and the downstream leak parameters are reconstructed analytically. The resulting realizations are again filtered through the spatial and hydraulic constraints, yielding probability distributions for the locations and emitter coefficients of both leaks.
\end{enumerate}

This constrained Monte Carlo framework enables an approximate localization of two simultaneous leaks using steady-state boundary measurements. Scaling the method beyond two leaks quickly increases the complexity. To detect 3 leaks, uniform sampling across 4 parameters is required, rather than 2. Additionally, there is no spatial bound between the second leak and the single-leak equivalent, as shown in \autoref{app:3_leak_eq_proof}. Based on these observations, it seems infeasible to extend the method beyond 2 leaks.

